%%%%%%%%%%%%%%%%%%%%%%%%%%%%%%%%%%%%%%%%%%%%%%%%%%%%%%%%%%%%%%%%%%%%%%%
%%  main.tex — demo del tema pablozd (beamer + metropolis)
%%  Compilar: latexmk -pdf main.tex   (usa biber para la bibliografía)
%%%%%%%%%%%%%%%%%%%%%%%%%%%%%%%%%%%%%%%%%%%%%%%%%%%%%%%%%%%%%%%%%%%%%%%
\documentclass[aspectratio=169,10pt]{beamer}

\usetheme{pablozd}          % tema (opcional: [banner=otra-imagen.png])
\usepackage{zd-paquetes}    % preámbulo de paquetes ordenado

\title[MDO]{Marcación Diferencial de Objetos}
\subtitle{Un template \themename para presentaciones de lingüística}
\author[P. Zdrojewski]{Pablo D. Zdrojewski}
\institute{Facultad de Filosofía y Letras · Universidad de Buenos Aires}
\date{6 de julio de 2026}
% Logos institucionales. El logo oficial de FILO:UBA se descarga una vez:
%   wget -O logo-filo-uba.png https://www.filo.uba.ar/themes/contrib/fedro-theme/logo.png
% (verificá que no sea la variante blanca para fondo oscuro; si lo es,
%  usá la versión en color del kit de identidad de la Facultad).
% Si el archivo existe se usa; si no, la demo cae en los logos de prueba.
\IfFileExists{logo-filo-uba.png}{%
  \titlegraphic{\zdlogo{logo-filo-uba.png}}%
}{%
  \titlegraphic{\zdlogo{logo-test1.png}\zdlogo{logo-test2.png}}%
}

\begin{document}

% ---------- Portada con banner ----------
\begin{frame}[plain]
  \titlepage
\end{frame}

% ---------- Índice ----------
\begin{frame}{Contenidos}
  \tableofcontents
\end{frame}

% =======================================================
\section{Ejemplos glosados}
% =======================================================

\begin{frame}{Ejemplos con \texttt{langsci-gb4e} y \texttt{leipzig}}
  El contraste básico de la MDO en español \citep{saab2021nonexistence}:
  \begin{exe}
    \ex\label{ex:mdo}
    \begin{xlist}
      \ex[]{
        \gll Juan vio *(a) María.\\
        Juan vio.\Third{}\Sg{} \Dom{} María\\
        \glt `Juan vio a María.'
      }
      \ex[*]{
        \gll Juan la\idx{i} vio a la casa\idx{i}.\\
        Juan \Cl{}.\Acc{} vio \Dom{} la casa\\
        \glt Duplicación de clíticos con objetos inanimados
      }
    \end{xlist}
  \end{exe}
  El ejemplo (\ref{ex:mdo}) muestra la interacción entre \alert{animacidad}
  y \alert{especificidad}.
\end{frame}

% =======================================================
\section{Árboles}
% =======================================================

\begin{frame}{Árboles con \texttt{forest}}
  \centering
  \begin{forest} arbolzd
    [S
      [SD [Juan]]
      [Sv
        [v\('\)
          [v]
          [SV
            [V [vio]]
            [SK
              [K [a]]
              [SD [María]]]]]]]
  \end{forest}
\end{frame}

% =======================================================
\section{Semántica formal}
% =======================================================

\begin{frame}{Bloques y notación semántica}
  \begin{block}{Denotación}
    $\llbracket \text{ver} \rrbracket^{g} =
      \lambda x \lambda e\,.\,\mathrm{ver}(e) \wedge \mathrm{Tema}(e, x)$
  \end{block}
  \begin{exampleblock}{Generalización}
    Todo objeto \([+\text{animado}, +\text{específico}]\) aparece marcado con \emph{a}.
  \end{exampleblock}
  \begin{alertblock}{Problema abierto}
    ¿La MDO es un fenómeno de Caso o de estructura informativa?
  \end{alertblock}
\end{frame}

\begin{frame}{Tablas con \texttt{booktabs} y columnas}
  \begin{multicols}{2}
    \begin{table}
      \centering
      \begin{tabular}{lcc}
        \toprule
        Objeto & Animado & MDO \\
        \midrule
        María        & sí & obligatoria \\
        el perro     & sí & variable \\
        la casa      & no & ausente \\
        \bottomrule
      \end{tabular}
    \end{table}
    \columnbreak
    \begin{itemize}
      \item Nodo principal
      \item Otro nodo
        \begin{itemize}
          \item Subnodo
            \begin{itemize}
              \item Sub-subnodo
            \end{itemize}
        \end{itemize}
    \end{itemize}
  \end{multicols}
\end{frame}

% ---------- Standout ----------
% [plain] oculta el pie de navegación en frames de cierre
\begin{frame}[standout, plain]
  ¿Preguntas?
\end{frame}

% ---------- Bibliografía ----------
\appendix
\nocite{*}
\begin{frame}[allowframebreaks]{Referencias}
  \printbibliography[heading=none]
\end{frame}

\end{document}
